\chapter[Weak Equation Contract]{The Weak Equation as an Elaboration Contract}
\label{chap:weak-equation-elaboration-contract}

The weak solution is the compiler's preferred shape for a variational statement. A strong equation asks
for an equality at every point where the expression is defined. A weak equation asks for vanishing against
all admissible tests. The second request may sound less exact, but for the compiler it is more honest: it
names the space of tests, the norm in which they are controlled, and the linear functional whose vanishing
must be checked. Nothing is lost by refusing pointwise speech when the state does not have pointwise
regularity. The exact statement is the one that matches the space in which the state lives.

Let \(\mathcal A\) be the activity functional. Its first Frechet variation at \(u\) in the direction
\(v\) is
\[
  D\mathcal A(u)[v]
    =
  \lim_{\epsilon\to 0}
    \frac{\mathcal A(u+\epsilon v)-\mathcal A(u)}{\epsilon},
\]
when the limit exists in the appropriate Sobolev sense. The principle of least activity becomes the weak
equation
\[
  D\mathcal A(u)[v] = 0
  \qquad \text{for all } v\in V .
\]
The compiler reads the universal quantifier as a contract generator. Each admissible test produces one
elaboration obligation. The derivative is not merely a calculus symbol; it is the linearized dependency of
the activity on a perturbation of the state, exposed in the form the checker can apply to a test term.

\section{The first variation as sensitivity}
\label{se:first-variation-sensitivity}

Frechet differentiation is the right derivative for the compiler because it respects the ambient space. A
partial derivative changes one coordinate while pretending the others are silent. A Frechet derivative
changes the state by an element of the same function space in which the state is being judged. That is the
compiler's discipline in analytic form: perturbations must have the type the equation says they have, and
the response to a perturbation must be linear and bounded in the norm named by the problem.

The derivative is therefore a sensitivity map,
\[
  D\mathcal A(u) : V \to \mathbb R,
\]
or into the scalar field carried by the construction. To prove that \(u\) is stationary is to prove that
this map is the zero functional. In strong language one might try to simplify the integrand until a
pointwise Euler--Lagrange equation appears. In weak language the compiler receives the thing it can
actually test: a functional waiting for inputs. If every accepted input reads zero, the stationary
condition is satisfied at the promised level.

\section{Least activity as zero visible descent}
\label{se:least-activity-zero-visible-descent}

The phrase ``least activity'' should not be read as an instruction to search an unbounded landscape until a
minimum is guessed. At the weak level it says that no admissible first-order perturbation lowers the
activity. The compiler can express that as zero visible descent: every test direction the contract admits
has first variation zero. When an inequality or convexity argument is also present, this stationary point
may be upgraded to a minimizer. Without that extra structure the compiler should not pretend to have more
than stationarity.

This separation is one of the main benefits of the weak contract. It tanges stationarity from minimization.
The first Frechet variation supplies the stationarity obligation. A second variation, coercivity estimate,
or convexity law supplies the additional minimization obligation. If those facts are present, they can be
checked as separate rows. If they are absent, the ledger remains honest: the first variation vanished on
the tested space, and no stronger result was smuggled in under the word ``least.''

\section{Quantifiers and cost}
\label{se:weak-quantifiers-and-cost}

The full weak equation quantifies over an infinite test space. The compiler cannot discharge that infinity
by enumeration. It needs either a symbolic argument that closes the universal, or a finite subspace whose
scope is declared. Those are different results. A universal weak theorem says \(D\mathcal A(u)[v]=0\) for
arbitrary \(v\) satisfying the type of admissibility. A Galerkin theorem says the same equality holds for
the basis tests, and therefore for their finite span. Both are exact when stated correctly. Confusion
begins only when the finite theorem is spoken as though it were already the universal one.

Cost records the difference. A universal proof spends effort on abstraction: bind an arbitrary test, use
the laws, discharge the expression without knowing the coefficients in advance. A Galerkin computation
spends effort on rows: choose a basis, assemble the finite residual, compute the entries, and close the
system. The compiler can host both modes, but the ledger has to name which mode is being used. The weak
equation as an elaboration contract keeps that naming visible.
